\section{Relations Over Types\label{sec:RelationsOnTypes}}

This section defines the following relations over types and operators:
\begin{itemize}
  \item Subtype (\TypingRuleRef{Subtype})
  \item Subtype Satisfaction (\TypingRuleRef{SubtypeSatisfaction})
  \item Type Satisfaction (\TypingRuleRef{TypeSatisfaction})
  \item The Lowest Common Ancestor of two types (\TypingRuleRef{LowestCommonAncestor})
  \item Applying a unary operator to a type (\TypingRuleRef{ApplyUnopType})
  \item Applying a binary operator to a pair of types (\TypingRuleRef{ApplyBinopTypes})
\end{itemize}

\TypingRuleDef{Subtype}
\hypertarget{def-supertypeterm}{}
The \emph{\subtypeterm} relation is a partial order over \underline{named types}.
The \emph{\supertypeterm} is the inverse relation.
That is, \tty\ is a \supertypeterm{} of \tsy\ if and only if \tsy\ is a \subtypeterm{} of \tty.

\ExampleDef{Subtypes and Supertypes}
The following table determines whether the type \vAbf{}
subtypes the type \vBbf{} with respect to the types
declared in \listingref{subtype}:\\
\begin{tabular}{llll}
  \textbf{type A} & \textbf{type B}   & \textbf{subtypes?}  & \textbf{reason}\\
\hline
  \texttt{subInt}     & \texttt{subInt}   & yes             & subtyping is reflexive for \namedtypesterm{}\\
  \texttt{subInt}     & \texttt{superInt} & yes             & declared as a \subtypeterm{}\\
  \texttt{superInt}   & \texttt{subInt}   & no              & subtyping is antisymmetric\\
  \texttt{subsubInt}  & \texttt{superInt} & yes             & subtyping is transitive\\
  \texttt{otherInt}   & \texttt{superInt} & no              & no chain of subtyping between the types\\
  \texttt{superInt}   & \texttt{integer}  & no              & \texttt{integer} is not a \namedtypeterm{}\\
  \texttt{integer}    & \texttt{integer}  & no              & \texttt{integer} is not a \namedtypeterm{}\\
\end{tabular}

\ASLListing{Subtypes and Supertypes}{subtype}{\typingtests/TypingRule.Subtype.asl}

\RenderRelation{is_subtype}

\RenderProseAndFormally{is_subtype}
\CodeSubsection{\SubtypeBegin}{\SubtypeEnd}{../types.ml}

\identr{NXRX} \identi{KGKS} \identi{MTML} \identi{JVRM} \identi{CHMP}

\TypingRuleDef{SubtypeSatisfaction}
\RenderRelation{subtype_satisfies}

This function assumes that both $\vt$ and $\vs$ are well-typed according to \chapref{Types}.

\ExampleDef{Subtype Satisfaction}
\listingref{subtypesat1} shows examples
where the types of the \rhsexpressions{} \subtypesatisfyterm{}
the types of the left-hand-side expressions.
\pagebreak
\ASLListing{Subtype Satisfaction}{subtypesat1}{\typingtests/TypingRule.SubtypeSatisfaction1.asl}

\listingref{subtypesat-bad1} and \listingref{subtypesat-bad2}
shows examples of nuanced \typingerrorsterm.
Specifically where a type consisting of a range of values does not \subtypesatisfyterm{}
a type consisting of one variable expression.
\ASLListing{Subtype satisfaction error 1}{subtypesat-bad1}{\typingtests/TypingRule.SubtypeSatisfaction.bad1.asl}
\ASLListing{Subtype satisfaction error 2}{subtypesat-bad2}{\typingtests/TypingRule.SubtypeSatisfaction.bad2.asl}

\listingref{subtypesat2} shows examples of legal and illegal assignments involving
subtyping.
\ASLListing{More Examples of subtype satisfaction}{subtypesat2}{\typingtests/TypingRule.SubtypeSatisfaction2.asl}

\listingref{subtypesat3} shows more examples of legal and illegal assignments involving
subtyping.
\ASLListing{Even more examples of subtype satisfaction}{subtypesat3}{\typingtests/TypingRule.SubtypeSatisfaction3.asl}

We employ the following helper function.
\RenderRelation{field_type}
This helper function is defined as follows:
\[
  \fieldtype(\fields, \id) \triangleq
  \begin{cases}
  \some{\vt}  & \text{ if }\{ \vtp \;|\; (\id,\vtp) \in \fields\} = \{\vt\}\\
  \None & \text{ otherwise}
  \end{cases}
\]


\RenderProseAndFormally{subtype_satisfies}
\identd{TRVR} \identi{SJDC} \identi{MHYB} \identi{TWTZ} \identi{GYSK} \identi{KXSD} \identi{KNXJ}
\CodeSubsection{\SubtypeSatisfactionBegin}{\SubtypeSatisfactionEnd}{../types.ml}

\TypingRuleDef{TypeSatisfaction}
\identr{FMXK}
\RenderRelation{type_satisfies}

The function assumes that both $\vt$ and $\vs$ are well-typed according to \secref{Types}.

\ExampleDef{Type-satisfaction Examples}
In \listingref{typing-typesat1},
\texttt{var pair: pairT = (1, dataT1)} is legal since the right-hand-side has
anonymous, non-primitive type \texttt{(integer, T1)}.
\ASLListing{Type satisfaction example}{typing-typesat1}{\typingtests/TypingRule.TypeSatisfaction1.asl}

\ExampleDef{More Type-satisfaction Examples}
In \listingref{typing-typesat2},
\texttt{pair = (1, dataAsInt);} is legal since the right-hand-side has anonymous,
primitive type \texttt{(integer, integer)}.
\ASLListing{Type satisfaction example}{typing-typesat2}{\typingtests/TypingRule.TypeSatisfaction2.asl}

\ExampleDef{Failing Type-satisfaction}
In \listingref{typing-typesat3},
\texttt{pair = (1, dataT2);} is illegal since the right-hand-side has anonymous,
non-primitive type \texttt{(integer, T2)} which does not subtype-satisfy named
type \texttt{pairT}.
\ASLListing{Type satisfaction example}{typing-typesat3}{\typingtests/TypingRule.TypeSatisfaction3.asl}

The specification in \listingref{typing-typesat-bad1} is ill-typed,
since \verb|integer{0..N}| does not \typesatisfyterm{} \verb|integer{0..M}|.
\ASLListing{Failing type satisfaction example}{typing-typesat-bad1}{\typingtests/TypingRule.TypeSatisfaction.bad1.asl}

\RenderProseAndFormally{type_satisfies}
\CodeSubsection{\TypeSatisfactionBegin}{\TypeSatisfactionEnd}{../types.ml}

\TypingRuleDef{CheckTypeSatisfaction}
\RenderRelation{check_type_satisfies}

The function assumes that both $\vt$ and $\vs$ are well-typed according to \secref{Types}.

\ExampleDef{Checking Type Satisfaction}
In \listingref{typing-typesat1},
checking whether \verb|(integer, T1)| \typesatisfiesterm{} \verb|pairT|
for the assignment \verb|var pair: pairT = (1, dataT1)| yields $\True$.

In \listingref{typing-typesat3}, checking whether \verb|(intege, T2)|
\typesatisfiesterm{} \verb|pairT| for the assignment \verb|pair = (1, dataT2);|
yields a \typingerrorterm.


\RenderProseAndFormally{check_type_satisfies}

\TypingRuleDef{LowestCommonAncestor}
Annotating a conditional expression (see \TypingRuleRef{ECond}),
requires finding a single type that can be used to annotate the results of both subexpressions.
We refer to such a type as a \emph{\Proselca}, or LCA, for short, and define it next.

\RenderRelation{lowest_common_ancestor}

\ExampleDef{Lowest Common Ancestor Examples}
\listingref{example-lca} shows examples of conditional expressions and the resulting
\emph{\Proselca}.
\pagebreak
\ASLListing{Lowest common ancestor}{example-lca}{\typingtests/TypingRule.LowestCommonAncestor.asl}

\listingref{example-lca2} shows an example of a lowest common ancestor of two \bitvectortypesterm{}
of different widths.
\ASLListing{Lowest common ancestor of two bitvectors}{example-lca2}{\typingtests/TypingRule.LowestCommonAncestor2.asl}

Note that, since we do not impose a canonical representation on types
(e.g., \verb|integer {1, 2}| is \typeequivalentterm{} to \verb|integer {1..2}|),
the lowest common ancestor is not unique.
We define $\lca(\tenv, \vt, \vs)$ to be any type $\vtp$ that is \typeequivalentterm{} to the lowest common ancestor of $\vt$ and $\vs$.

\RenderProseAndFormally{lowest_common_ancestor}
\CodeSubsection{\LowestCommonAncestorBegin}{\LowestCommonAncestorEnd}{../types.ml}
\identr{YZHM}

\TypingRuleDef{ApplyUnopType}
\RenderRelation{apply_unop_type}

\ExampleDef{Applying Unary Operations to Types}
\listingref{apply-unop-type} shows examples of typing applications of unary operations.
\ASLListing{Applying unary operations to types}{apply-unop-type}{\typingtests/TypingRule.ApplyUnopType.asl}


\RenderProseAndFormally{apply_unop_type}
\CodeSubsection{\ApplyUnopTypeBegin}{\ApplyUnopTypeEnd}{../Typing.ml}

\TypingRuleDef{NegateConstraint}
\RenderRelation{negate_constraint}

\ExampleRef{Applying Unary Operations to Types} shows examples of negating single expression
constraints and range constraints.


\RenderProseAndFormally{negate_constraint}

\TypingRuleDef{ApplyBinopTypes}
\RenderRelation{apply_binop_types}

\ExampleDef{Filtering Constraints for Binary Operations}
Since binary operations may dynamically fail, the typechecker may remove values
that will definitely lead to \dynamicerrorsterm.
For example, in \listingref{apply-binop-constraints2},
the constraint \verb|integer{-1..1}| serves for the denominator in \verb|x DIV y|,
and since \verb|-1| and \verb|0| will always lead to a \dynamicerrorterm,
the typechecker removes them for the purpose of typing \verb|x DIV y|,
which leaves only \verb|1| and the final type for \verb|z| is therefore
\verb|integer{2, 4}|.

\ASLListing{Removing values from constraints}{apply-binop-constraints2}
{\typingtests/TypingRule.ApplyBinopTypes.constraints2.asl}

\ExampleDef{Applying Binary Operations to Types}
\listingref{apply-binop-type} shows examples of typing binary operations.
\ASLListing{Applying binary operations to types}{apply-binop-type}
          {\typingtests/TypingRule.ApplyBinopTypes.asl}

\ExampleDef{Applying Binary Operations to Constrained Integers}
\listingref{apply-binop-constraints} shows examples of typing binary operations
applied to \constrainedintegerterm{} types.

Importantly, note that there is no AST for applying binary operations on
constraints.
For example, given a range constraint \verb|A..B| and an exact constraint |2|,
there is no AST to express \verb|(A..B) * 2|. Therefore, the constraints for typing
\verb|ab_times_2| approximate the set of values for \verb|(A..B) * 2| via four
range constraints. More precisely, they are a superset of the values for \verb|(A..B) * 2|.

\pagebreak
\ASLListing{Applying binary operations to constrained integers}{apply-binop-constraints}
          {\typingtests/TypingRule.ApplyBinopTypes.constraints.asl}

\RenderProseAndFormally{apply_binop_types}

The signatures for the arithmetic operations for the \realtypeterm{} are as follows:
\RenderConstant{arith_real_signatures}

\CodeSubsection{\applybinoptypesBegin}{\applybinoptypesEnd}{../Typing.ml}

\identr{BKNT} \identr{ZYWY} \identr{BZKW}
\identr{KFYS} \identr{KXMR} \identr{SQXN} \identr{MRHT} \identr{JGWF}
\identr{TTGQ} \identi{YHML} \identi{YHRP} \identi{VMZF} \identi{YXSY}
\identi{LGHJ} \identi{RXLG}

\TypingRuleDef{FindNamedLCA}
\RenderRelation{named_lowest_common_ancestor}
\RenderRelation{supers}
\[
  \supers(\tenv, \vt) \triangleq
  \begin{cases}
    \{\vt\} \cup \supers(\vs) & \text{ if }\tenv.\staticenvsG.\subtypes(\vt) = \vs\\
    \{\vt\}  & \text{ otherwise } (\text{that is, }\tenv.\staticenvsG.\subtypes(\vt) = \bot)\\
  \end{cases}
\]

\ExampleDef{Finding Named Lowest Common Ancestors}
In \listingref{find-named-lca},
the set of named supertypes for \verb|B2| is \verb|{A1, A2, B1, B2}|,
set of named supertypes for \verb|C2| is \verb|{A1, A2, C1, C2}|,
therefore the named lowest common ancestor of \verb|B2| and \verb|C2|
is \verb|A2|, while \verb|B2| and \verb|D1| have no
named lowest common ancestor.

\ASLListing{Finding named lowest common ancestors}{find-named-lca}{\typingtests/TypingRule.FindNamedLCA.asl}


\RenderProseAndFormally{named_lowest_common_ancestor}

\TypingRuleDef{AnnotateConstraintBinop}
\RenderRelation{annotate_constraint_binop}

The operator $\op$ is assumed to be only one of the operators in the following set:
$\{\SHL, \SHR, \POW, \MOD, \DIVRM, \SUB, \MUL, \ADD, \DIV\}$.
The rule employs \\ $\binopisexploding$ to decide whether range constraints can be maintained
as range constraints or have to be converted to a list of exact constraints.

\ExampleDef{Annotating Constraints for Binary Operations}
Applying $\ADD$ to
$\{\AbbrevConstraintRange{\ELInt{2}}{\ELInt{4}}\}$ and
$\{\AbbrevConstraintExact{\ELInt{2}}\}$ \\ results in
$\{\AbbrevConstraintRange{\ELInt{4}}{\ELInt{6}}\}$,\\
since $\binopisexploding(\ADD) \typearrow \False$
while applying $\MUL$ to the same lists of constraints results in\\
$\{\AbbrevConstraintExact{\ELInt{4}}, \AbbrevConstraintExact{\ELInt{6}}, \AbbrevConstraintExact{\ELInt{8}}\}$,\\
since $\binopisexploding(\MUL) \typearrow \True$.

Annotating the constraints involves applying symbolic reasoning and in particular filtering out values that
will definitely result in a \dynamicerrorterm{}.

Also see \ExampleRef{Applying Binary Operations to Constrained Integers}.

\RenderProseAndFormally{annotate_constraint_binop}
\CodeSubsection{\AnnotateConstraintBinopBegin}{\AnnotateConstraintBinopEnd}{../StaticOperations.ml}

\TypingRuleDef{BinopFilterRhs}
\RenderRelation{binop_filter_rhs}

\ExampleDef{Filtering Right-hand-side Constraints}
An example of filtering constraints appears in \listingref{apply-binop-constraints},
where the constraints inferred for \verb|a_div| filter out \verb|-5..0|
from the constraint \verb|-5..3|, thus avoiding including constraints
\verb|A DIV -5| and \verb|A DIV 0|.

\ProseParagraph
\OneApplies
\begin{itemize}
  \item \AllApplyCase{greater\_or\_equal}
  \begin{itemize}
    \item $\op$ is one of $\SHL$, $\SHR$, and $\POW$;
    \item define $\vf$ as the specialization of $\refineconstraintbysign$ for the predicate
          $\lambda x.\ x \geq 0$, which is $\True$ if and only if the tested number is greater than or equal to $0$;
    \item refining the list of constraints $\cs$ with $\vf$ and $\vapprox$ via $\refineconstraints$ yields \\
          $\newcs$\ProseTerminateAs{\CannotUnderapproximate, \CannotOverapproximate};
    \item checking whether $\newcs$ is empty yields $\True$\ProseTerminateAs{\BadOperands}.
  \end{itemize}

  \item \AllApplyCase{greater\_than}
  \begin{itemize}
    \item $\op$ is one of $\MOD$, $\DIV$, and $\DIVRM$;
    \item define $\vf$ as the specialization of $\refineconstraintbysign$ for the predicate
          $\lambda x.\ x > 0$, which is $\True$ if and only if the tested number is greater than $0$;
    \item refining the list of constraints $\cs$ with $\vf$ and $\vapprox$ via $\refineconstraints$ yields \\
          $\newcs$\ProseTerminateAs{\CannotUnderapproximate, \CannotOverapproximate};
    \item checking whether $\newcs$ is empty yields $\True$\ProseTerminateAs{\BadOperands}.
  \end{itemize}

  \item \AllApplyCase{no\_filtering}
  \begin{itemize}
    \item $\op$ is one of $\SUB$, $\MUL$, and $\ADD$;
    \item $\newcs$ is $\cs$.
  \end{itemize}
\end{itemize}

\FormallyParagraph
\begin{mathpar}
\inferrule[greater\_or\_equal]{
  \op \in \{\SHL, \SHR, \POW\}\hva\\\\
  \refineconstraintbysign(\tenv, \lambda x.\ x \geq 0) \typearrow \vf\\
  \refineconstraints(\vapprox, \cs, \vf) \typearrow \newcs \terminateas \CannotUnderapproximate, \CannotOverapproximate\hva\\\\
  \techeck(\newcs \neq \emptylist, \BadOperands) \typearrow \True\OrTypeError
}{
  \binopfilterrhs(\vapprox, \tenv, \op, \cs) \typearrow \newcs
}
\end{mathpar}

\begin{mathpar}
\inferrule[greater\_than]{
  \op \in \{\MOD, \DIV, \DIVRM\}\hva\\\\
  \refineconstraintbysign(\tenv, \lambda x.\ x > 0) \typearrow \vf\\
  \refineconstraints(\vapprox, \cs, \vf) \typearrow \newcs \terminateas \CannotUnderapproximate, \CannotOverapproximate\hva\\\\
  \techeck(\newcs \neq \emptylist, \BadOperands) \typearrow \True\OrTypeError
}{
  \binopfilterrhs(\vapprox, \tenv, \op, \cs) \typearrow \newcs
}
\end{mathpar}

\begin{mathpar}
\inferrule[no\_filtering]{
  \op \in \{\SUB, \MUL, \ADD\}
}{
  \binopfilterrhs(\vapprox, \op, \cs) \typearrow \overname{\cs}{\newcs}
}
\end{mathpar}

\TypingRuleDef{RefineConstraintBySign}
\RenderRelation{refine_constraint_by_sign}

\ExampleDef{Refining Constraints with a Sign Predicate}
The specification in \listingref{RefineConstraintBySign}
shows examples of how the constraints for \verb|y|
are refined for the \DIV{} operator via the sign predicate
that tests whether an integer is positive: $\lambda x.\ x > 0$.

\ASLListing{Refining Constraints with a Sign Predicate}{RefineConstraintBySign}{\typingtests/TypingRule.RefineConstraintBySign.asl}

\ProseParagraph
\OneApplies
\begin{itemize}
  \item \AllApplyCase{exact\_reduces\_to\_z}
  \begin{itemize}
    \item $\vc$ is an exact constraint for the expression $\ve$, that is, $\ConstraintExact(\ve)$;
    \item applying $\reducetozopt$ to $\ve$ in $\tenv$, in order to symbolically simplify $\ve$ to an integer,
          yields $\some{\vz}$;
    \item $\vcopt$ is $\some{\vc}$ if $\vp$ holds for $\vz$ and $\None$ otherwise.
  \end{itemize}

  \item \AllApplyCase{exact\_does\_not\_reduce\_to\_z}
  \begin{itemize}
    \item $\vc$ is an exact constraint for the expression $\ve$, that is, $\ConstraintExact(\ve)$;
    \item applying $\reducetozopt$ to $\ve$ in $\tenv$, in order to symbolically simplify $\ve$ to an integer,
          yields $\None$;
    \item $\vcopt$ is $\some{\vc}$.
  \end{itemize}

  \item \AllApplyCase{range\_both\_reduce\_to\_z}
  \begin{itemize}
    \item $\vc$ is a range constraint for the expressions $\veone$ and $\vetwo$, that is, \\
          $\ConstraintRange(\veone, \vetwo)$;
    \item applying $\reducetozopt$ to $\veone$ in $\tenv$, in order to symbolically simplify $\veone$ to an integer,
          yields $\some{\vzone}$;
    \item applying $\reducetozopt$ to $\vetwo$ in $\tenv$, in order to symbolically simplify $\vetwo$ to an integer,
          yields $\some{\vztwo}$;
    \item \OneApplies{} (defining $\vcopt$)
    \begin{itemize}
      \item if $\vp$ is $\True$ for both $\vzone$ and $\vztwo$, define $\vcopt$ as $\some{\vc}$;
      \item if $\vp$ is $\False$ for $\vzone$ and $\True$ for $\vztwo$, define $\vcopt$ as the optional range constraint
            where the bottom expression is the literal expression for $0$ if $\vp$ holds for $0$ and the literal expression for $1$ otherwise,
            and the top expression is $\vetwo$;
      \item if $\vp$ is $\True$ for $\vzone$ and $\False$ for $\vztwo$, define $\vcopt$ as the optional range constraint
            where the bottom expression is $\veone$ and the top expression is the literal expression for $0$ if $\vp$ holds for $0$
            and the literal expression for $-1$ otherwise;
      \item if $\vp$ is $\False$ for both $\vzone$ and $\vztwo$, define $\vcopt$ as $\None$.
    \end{itemize}
  \end{itemize}

  \item \AllApplyCase{only\_e1\_reduces\_to\_z}
  \begin{itemize}
    \item $\vc$ is a range constraint for the expressions $\veone$ and $\vetwo$, that is, \\
          $\ConstraintRange(\veone, \vetwo)$;
    \item applying $\reducetozopt$ to $\veone$ in $\tenv$, in order to symbolically simplify $\veone$ to an integer,
          yields $\some{\vzone}$;
    \item applying $\reducetozopt$ to $\vetwo$ in $\tenv$, in order to symbolically simplify $\vetwo$ to an integer,
          yields $\None$;
    \item \OneApplies{} (defining $\vcopt$):
    \begin{itemize}
      \item if $\vp$ is $\True$ for $\vzone$, define $\vcopt$ as $\some{\vc}$;
      \item if $\vp$ is $\False$ for $\vzone$, define $\vcopt$ as the optional range constraint with the bottom expression
            as the literal expression for $0$ if $\vp$ holds for $0$ and the literal expression for $1$ otherwise,
            and the top expression $\vetwo$.
    \end{itemize}
  \end{itemize}

  \item \AllApplyCase{only\_e2\_reduces\_to\_z}
  \begin{itemize}
    \item $\vc$ is a range constraint for the expressions $\veone$ and $\vetwo$, that is, \\
          $\ConstraintRange(\veone, \vetwo)$;
    \item applying $\reducetozopt$ to $\veone$ in $\tenv$, in order to symbolically simplify $\veone$ to an integer,
          yields $\None$;
    \item applying $\reducetozopt$ to $\vetwo$ in $\tenv$, in order to symbolically simplify $\vetwo$ to an integer,
          yields $\some{\vztwo}$;
    \item One of the following applies (defining $\vcopt$):
    \begin{itemize}
      \item if $\vp$ is $\True$ for $\vztwo$, define $\vcopt$ as $\some{\vc}$;
      \item if $\vp$ is $\False$ for $\vztwo$, define $\vcopt$ as the optional range constraint with the bottom expression
            $\veone$ and the top expression the literal expression for $0$ if $\vp$ holds for $0$ and the literal expression for $-1$ otherwise.
    \end{itemize}
  \end{itemize}

  \item \AllApplyCase{none\_reduce\_to\_z}
  \begin{itemize}
    \item $\vc$ is a range constraint for the expressions $\veone$ and $\vetwo$, that is, \\
          $\ConstraintRange(\veone, \vetwo)$;
    \item applying $\reducetozopt$ to $\veone$ in $\tenv$, in order to symbolically simplify $\veone$ to an integer,
          yields $\None$;
    \item applying $\reducetozopt$ to $\vetwo$ in $\tenv$, in order to symbolically simplify $\vetwo$ to an integer,
          yields $\None$;
    \item \Proseeqdef{$\vcopt$}{$\vc$}.
  \end{itemize}
\end{itemize}

\FormallyParagraph
\begin{mathpar}
\inferrule[exact\_reduces\_to\_z]{
  \reducetozopt(\tenv, \ve) \typearrow \some{\vz}\\
  \vcopt \eqdef \ifthenelseop[H]{\vp(\vz)}{\some{\vc}}{\None}
}{
  \refineconstraintbysign(\tenv, \vp, \overname{\ConstraintExact(\ve)}{\vc}) \typearrow \vcopt
}
\end{mathpar}

\begin{mathpar}
\inferrule[exact\_does\_not\_reduce\_to\_z]{
  \reducetozopt(\tenv, \ve) \typearrow \None
}{
  \refineconstraintbysign(\tenv, \vp, \overname{\ConstraintExact(\ve)}{\vc}) \typearrow \overname{\some{\vc}}{\vcopt}
}
\end{mathpar}

\begin{mathpar}
\inferrule[range\_both\_reduce\_to\_z]{
  \reducetozopt(\tenv, \veone) \typearrow \some{\vzone}\\
  \reducetozopt(\tenv, \vetwo) \typearrow \some{\vztwo}\\
  \vzeroorone \eqdef \ifthenelseop[H]{\vp(0)}{\ELInt{0}}{\ELInt{1}}\hva\\
  \vzeroorminusone \eqdef \ifthenelseop[H]{\vp(0)}{\ELInt{0}}{\ELInt{-1}}\\
  \vrangezeroorone \eqdef \ConstraintRange(\vzeroorone, \vetwo)\hva\\
  \vrangezeroorminusone \eqdef \ConstraintRange(\veone, \vzeroorminusone)\\
  {
  \vcopt \eqdef \begin{cases}
    \some{\vc}& \text{if }\vp(\vzone) \land \vp(\vztwo)\\
    \some{\vrangezeroorone}& \text{if }\neg\vp(\vzone) \land \vp(\vztwo)\\
    \some{\vrangezeroorminusone}& \text{if }\vp(\vzone) \land \neg\vp(\vztwo)\\
    \None& \text{if }\neg\vp(\vzone) \land \neg\vp(\vztwo)\\
  \end{cases}
  }
}{
  \refineconstraintbysign(\tenv, \vp, \overname{\ConstraintRange(\veone, \vetwo)}{\vc}) \typearrow \vcopt
}
\end{mathpar}

\begin{mathpar}
\inferrule[only\_e1\_reduces\_to\_z]{
  \vc = \ConstraintRange(\veone, \vetwo)\\
  \reducetozopt(\tenv, \veone) \typearrow \some{\vzone}\\
  \reducetozopt(\tenv, \vetwo) \typearrow \None\hva\\
  \vzeroorone \eqdef \ifthenelseop[H]{\vp(0)}{\ELInt{0}}{\ELInt{1}}\\
  \vrangezeroorone \eqdef \ConstraintRange(\vzeroorone, \vetwo)\\
  {
  \vcopt \eqdef  \begin{cases}
    \some{\vc}& \text{if }\vp(\vzone)\\
    \some{\vrangezeroorone}& \text{otherwise}\\
  \end{cases}
  }
}{
  \refineconstraintbysign(\tenv, \vp, \vc) \typearrow \vcopt
}
\end{mathpar}

\begin{mathpar}
\inferrule[only\_e2\_reduces\_to\_z]{
  \vc = \ConstraintRange(\veone, \vetwo)\\
  \reducetozopt(\tenv, \veone) \typearrow \None\\
  \reducetozopt(\tenv, \vetwo) \typearrow \some{\vztwo}\\
  \vzeroorminusone \eqdef \ifthenelseop[H]{\vp(0)}{\ELInt{0}}{\ELInt{-1}}\hva\\
  \vrangezeroorminusone \eqdef \ConstraintRange(\veone, \vzeroorminusone)\\
  {
  \vcopt \eqdef \begin{cases}
    \some{\vc}& \text{if }\vp(\vztwo)\\
    \some{\vrangezeroorminusone}& \text{otherwise}\\
  \end{cases}
  }
}{
  \refineconstraintbysign(\tenv, \vp, \vc) \typearrow \vcopt
}
\end{mathpar}

\begin{mathpar}
\inferrule[none\_reduce\_to\_z]{
  \vc = \ConstraintRange(\veone, \vetwo)\\
  \reducetozopt(\tenv, \veone) \typearrow \None\\
  \reducetozopt(\tenv, \vetwo) \typearrow \None
}{
  \refineconstraintbysign(\tenv, \vp, \vc) \typearrow \overname{\vc}{\vcopt}
}
\end{mathpar}

\TypingRuleDef{ReduceToZOpt}
\RenderRelation{reduce_to_z_opt}

\ExampleDef{Reducing Expressions to Optional Integers}
The specification in \listingref{ReduceToZOpt} shows two examples of expressions
where one can be reduced into an integer constant, which is then returned in an optional,
and one that cannot be reduced into an integer constant, which yields $\None$.
\ASLListing{Reducing expressions to optional integers}{ReduceToZOpt}{\typingtests/TypingRule.ReduceToZOpt.asl}


\RenderProseAndFormally{reduce_to_z_opt}

\TypingRuleDef{RefineConstraints}
\RenderRelation{refine_constraints}

\ExampleDef{Filtering Optional Constraints}
The specification in \listingref{RefineConstraintBySign}
shows how certain constraints are filtered out.
For example, $\AbbrevConstraintExact{\AbbrevEBinop{\DIV}{\AbbrevEVar{\vA}}{\ELInt{0}}}$.
Similarly, since both expressions in the constraint \verb|-4..-3| fail the sign predicate
used by $\refineconstraintbysign$, the resulting $\None$ is also filtered out.

\ProseParagraph
\OneApplies
\begin{itemize}
  \item \AllApplyCase{empty}
  \begin{itemize}
    \item $\cs$ is the empty list;
    \item $\newcs$ is the empty list.
  \end{itemize}

  \item \AllApply
  \begin{itemize}
    \item $\cs$ is the list with $\vc$ as its \head\ and $\csone$ as its \tail;
    \item applying $\vf$ to $\vc$ yields $\None$;
    \item applying $\refineconstraints$ to $\vapprox$, $\vf$ and $\csone$ yields $\csonep$;
    \item \OneApplies
    \begin{itemize}
      \item \AllApplyCase{non\_empty\_none\_precise}
      \begin{itemize}
        \item $\csonep$ is not the empty list;
        \item $\newcs$ is $\csonep$.
      \end{itemize}

      \item \AllApplyCase{non\_empty\_none\_approx}
      \begin{itemize}
        \item $\csonep$ is not the empty list;
        \item the result is $\CannotOverapproximate$ if $\vapprox$ is $\Over$ and $\CannotUnderapproximate$, otherwise.
      \end{itemize}
    \end{itemize}
  \end{itemize}

  \item \AllApplyCase{non\_empty\_some}
  \begin{itemize}
    \item $\cs$ is the list with $\vc$ as its \head\ and $\csone$ as its \tail;
    \item applying $\vf$ to $\vc$ yields $\some{\vcp}$;
    \item applying $\refineconstraints$ to $\vapprox$, $\vf$ and $\csone$ yields $\csonep$;
    \item $\newcs$ is the list with $\vcp$ as its \head\ and $\csonep$ as its \tail.
  \end{itemize}
\end{itemize}

\FormallyParagraph
\begin{mathpar}
\inferrule[empty]{}{
  \refineconstraints(\vapprox, \tenv, \vf, \overname{\emptylist}{\cs}) \typearrow \overname{\emptylist}{\newcs}
}
\end{mathpar}

\begin{mathpar}
\inferrule[non\_empty\_none\_precise]{
  \vf(\vc) \typearrow \None\\
  \refineconstraints(\vapprox, \vf, \csone) \typearrow \csonep\\
  \csonep \neq \emptylist
}{
  \refineconstraints(\vapprox, \vf, \overname{[\vc]\concat \csone}{\cs}) \typearrow \overname{\csonep}{\newcs}
}
\end{mathpar}

\begin{mathpar}
\inferrule[non\_empty\_none\_approx]{
  \vf(\vc) \typearrow \None\\
  \refineconstraints(\vapprox, \vf, \csone) \typearrow \csonep\\
  \csonep = \emptylist
}{
  \refineconstraints(\vapprox, \vf, \overname{[\vc]\concat \csone}{\cs}) \typearrow \ifthenelseop[H]{\vapprox=\Over}{\CannotOverapproximate}{\CannotUnderapproximate}
}
\end{mathpar}

\begin{mathpar}
\inferrule[non\_empty\_some]{
  \vf(\vc) \typearrow \some{\vcp}\\
  \refineconstraints(\vapprox, \vf, \csone) \typearrow \csonep\\
}{
  \refineconstraints(\vapprox, \vf, \overname{[\vc]\concat \csone}{\cs}) \typearrow \overname{[\vcp] \concat \csonep}{\newcs}
}
\end{mathpar}

\TypingRuleDef{RefineConstraintForDiv}
\RenderRelation{refine_constraint_for_div}

See \ExampleRef{Filtering and Reducing Division-based Constraints}.


\RenderProseAndFormally{refine_constraint_for_div}
\CodeSubsection{\RefineConstraintForDIVBegin}{\RefineConstraintForDIVEnd}{../types.ml}

\TypingRuleDef{FilterReduceConstraintDiv}
\RenderRelation{filter_reduce_constraint_div}

\ExampleDef{Filtering and Reducing Division-based Constraints}
The specification in \listingref{FilterReduceConstraintDiv}
shows how different constraints formed by dividing the constraints of \verb|x| by \verb|2|
in order to annotate \verb|z|.
The resulting constraints are listed in the type annotation for \verb|z_typed|.
Specifically, notice the following:
\begin{itemize}
  \item the constraint for \verb|3| divided by \verb|2| is filtered away,
        since the result is not an integer,
  \item the constraint for \verb|7..10| divided by \verb|2|
        effectively rounds up $7/2$ into \verb|3|,
  \item the constraint for \verb|14..17| divided by \verb|2| effectively rounds down
        $17/2$ into \verb|16|, and
  \item the constraint for \verb|16..15| divided by \verb|2| is filtered away,
        since it is not a legal \rangeconstraintterm.
\end{itemize}
\ASLListing{Filtering and Reducing Division-based Constraints}{FilterReduceConstraintDiv}{\typingtests/TypingRule.FilterReduceConstraintDiv.asl}


\RenderProseAndFormally{filter_reduce_constraint_div}

\TypingRuleDef{GetLiteralDivOpt}
\RenderRelation{get_literal_div_opt}

\ExampleDef{Obtaining Division Expressions Optionally Resulting in Literals}
In \listingref{FilterReduceConstraintDiv},
the following are examples of $\getliteraldivopt$ applications that are used
for generating the constraints for \verb|z|:
\[
\begin{array}{rcl}
\getliteraldivopt(\AbbrevEBinop{\DIV}{\ELInt{14}}{\ELInt{2}}) &\typearrow& \some{(14, 2)}\\
\getliteraldivopt(\AbbrevEBinop{\DIV}{\ELInt{7}}{\ELInt{2}}) &\typearrow& \some{(7, 2)}\\
\getliteraldivopt(\AbbrevEBinop{\DIV}{\AbbrevEVar{\vA}}{\ELInt{2}}) &\typearrow& \None\\
\end{array}
\]

\RenderProseAndFormally{get_literal_div_opt}

\TypingRuleDef{MakeConstraintRangeOpt}
\RenderRelation{make_constraint_range_opt}

\ExampleDef{Constructing Optional Range Constraints}
The following are examples of $\makeconstraintrangeopt$ applications:
\[
\begin{array}{rcl}
\makeconstraintrangeopt(42, 42) &\typearrow& \some{\AbbrevConstraintExact{\ELInt{42}}}\\
\makeconstraintrangeopt(0, 42) &\typearrow& \some{\AbbrevConstraintRange{\ELInt{0}}{\ELInt{42}}}\\
\makeconstraintrangeopt(42, 0) &\typearrow& \None\\
\end{array}
\]

\RenderProseAndFormally{make_constraint_range_opt}

\TypingRuleDef{ExplodeIntervals}
\RenderRelation{explode_intervals}

\ExampleDef{Exploding Intervals}
In \listingref{ExplodeIntervals}, annotating the constraints for \verb|y|
involves exploding each of the constraints for \verb|{C..0, 0..A}|,
which does not incur a precision loss.
On the other hand, annotating the constraints for \verb|z|
involves exploding each of the constraints for \verb|C..0, 0..B|
where exploding the first constraint \verb|C..0| does not incur precision loss,
but exploding the second constraint \verb|0..B| does incur precision loss,
and so overall the result is that there is precision loss.
\ASLListing{Exploding intervals}{ExplodeIntervals}{\typingtests/TypingRule.ExplodeIntervals.bad.asl}


\RenderProseAndFormally{explode_intervals}

\TypingRuleDef{ExplodeConstraint}
\RenderRelation{explode_constraint}

See \ExampleRef{Intervals too Large to Enumerate}.


\RenderProseAndFormally{explode_constraint}

\TypingRuleDef{IntervalTooLarge}
\RenderRelation{interval_too_large}

\ExampleDef{Intervals too Large to Enumerate}
In \listingref{IntervalTooLarge}, the size of the interval used to annotate
\verb|z| is less than $\maxexplodedintervalsize$.
\ASLListing{An enumerated interval}{IntervalTooLarge}{\typingtests/TypingRule.IntervalTooLarge.asl}

In \listingref{IntervalTooLarge-bad}, the size of the interval used to annotate
\verb|y| is just below $\maxexplodedintervalsize$, but the size of the interval
used to annotate \verb|z| is equal to $\maxexplodedintervalsize$, which is too
large to be enumerated.
\ASLListing{An interval too large to enumerate}{IntervalTooLarge-bad}{\typingtests/TypingRule.IntervalTooLarge.bad.asl}

\RenderProseAndFormally{interval_too_large}

\TypingRuleDef{BinopIsExploding}
\RenderRelation{binop_is_exploding}

See \ExampleRef{Annotating Constraints for Binary Operations}.

\RenderProseAndFormally{binop_is_exploding}

\TypingRuleDef{BitFieldsIncluded}
\RenderRelation{bitfields_included}

\ExampleDef{Checking Whether all Bitfields are Included}
In \listingref{MemBfs},
all bitfields of the type \\
\verb|bits(8) {[0] flag, [7:1] data { [2:0] low }}|
are included in the type \\
\verb|FlaggedPacket|.

In \listingref{MemBfs-bad}, not all bitfields of the type \verb|bits(8) {[0] flag, [1] lsb}|
are included in the type \verb|FlaggedPacket|. Specifically, the bitfield \verb|lsb| is
not included in \verb|FlaggedPacket|.


\RenderProseAndFormally{bitfields_included}

\TypingRuleDef{MemBfs}
\RenderRelation{mem_bfs}

\ExampleDef{Checking Bitfield Membership}
The specification in \listingref{MemBfs} is well-typed.
Specifically, the bitfield \verb|flag| in the type
\verb|bits(8) {[0] flag, [7:1] data { [2:0] low }}|
is a member of the set of bitfields in the type \verb|FlaggedPacket|.
\ASLListing{No missing bitfields}{MemBfs}{\typingtests/TypingRule.MemBfs.asl}

In \listingref{MemBfs-bad}, the bitfield \verb|lsb| is missing from
the set of bitfields in the type \\ \verb|FlaggedPacket|, which is why the specification
is ill-typed.
\ASLListing{A missing bitfield}{MemBfs-bad}{\typingtests/TypingRule.MemBfs.bad.asl}


\RenderProseAndFormally{mem_bfs}

\TypingRuleDef{CheckStructure}
\RenderRelation{check_structure_label}

See \ExampleRef{The Structure of a Type}.


\RenderProseAndFormally{check_structure_label}

\TypingRuleDef{ToWellConstrained}
\RenderRelation{to_well_constrained}

\ExampleDef{Converting Parameterized Integer Types to Well-constrained Integer Types}
The following are two examples of applying $\towellconstrained$ to various types:
\begin{mathpar}
\inferrule{}{
  {
    \begin{array}{r}
  \towellconstrained(\overname{\TInt(\Parameterized(\vx))}{\vt}) \typearrow\\ \TInt(\WellConstrained(\ConstraintExact(\EVar(\vx))))
    \end{array}
  }
}
\end{mathpar}

\begin{mathpar}
\inferrule{}{
  {
    \begin{array}{r}
  \towellconstrained(\overname{\TInt(\unconstrainedinteger)}{\vt}) \typearrow\\ \TInt(\unconstrainedinteger)
    \end{array}
  }
}
\end{mathpar}


\RenderProseAndFormally{to_well_constrained}

\TypingRuleDef{GetWellConstrainedStructure}
\RenderRelation{get_well_constrained_structure}

\ExampleDef{Obtaining the Well-constrained Structure of a Type}
In \listingref{GetWellConstrainedStructure},
annotating the expression \verb|-N| involves obtaining the well-constrained structure of the
type of \verb|N|, which yields \\
$\TInt(\WellConstrained([\ConstraintExact(\EVar(\vN))]))$.
Similarly, the for expression \verb|-x|,
obtaining the well-constrained structure of the type for $\vx$ yields \\
$\unconstrainedinteger$.
\ASLListing{Obtaining the well constrained structure of a type}{GetWellConstrainedStructure}{\typingtests/TypingRule.GetWellConstrainedStructure.asl}


\RenderProseAndFormally{get_well_constrained_structure}

\TypingRuleDef{GetBitvectorWidth}
\RenderRelation{get_bitvector_width}

See \ExampleRef{Obtaining the Integer Corresponding to a Bitvector Width}.


\RenderProseAndFormally{get_bitvector_width}

\CodeSubsection{\GetBitvectorWidthBegin}{\GetBitvectorWidthEnd}{../Typing.ml}

\TypingRuleDef{GetBitvectorConstWidth}
\RenderRelation{get_bitvector_const_width}

\ExampleDef{Obtaining the Integer Corresponding to a Bitvector Width}
In \listingref{lesetfields}, annotating the \assignableexpression{} \verb|x.[status, time, data]|
requires obtaining the widths of the fields \verb|status|, \verb|time|, and \verb|data|
whose corresponding types are \verb|bit|, \verb|bits(16)|, and \verb|bits(8)|,
which yield $1$, $16$, and $8$, respectively.


\RenderProseAndFormally{get_bitvector_const_width}

\CodeSubsection{\GetBitvectorConstWidthBegin}{\GetBitvectorConstWidthEnd}{../Typing.ml}

\TypingRuleDef{CheckBitsEqualWidth}
\RenderRelation{check_bits_equal_width}

\ExampleDef{Ensuring That Two Bitvectors Have Equal Width}
In \listingref{CheckBitsEqualWidth}, annotating the bitfield \verb|info|
requires checking that its declared width given by \verb|bits(4)| is equal to
the width defined by the corresponding slice \verb|[6:3]|.

In addition, annotating the expression \verb|x AND y| requires checking that
the widths of \verb|x| and \verb|y| are equal.
\ASLListing{Ensuring that two bitvectors Have equal width}{CheckBitsEqualWidth}{\typingtests/TypingRule.CheckBitsEqualWidth.asl}

The specification in \listingref{CheckBitsEqualWidth-bad}
is ill-typed, since \verb|M| and \verb|N| are not necessarily equal, even though they have the same type.
\ASLListing{Two bitvectors of possibly different widths}{CheckBitsEqualWidth-bad}{\typingtests/TypingRule.CheckBitsEqualWidth.bad.asl}

The specification in \listingref{CheckBitsEqualWidth-bad2} is ill-typed for a similar reason.
\ASLListing{Two bitvectors of possibly different widths}{CheckBitsEqualWidth-bad2}{\typingtests/TypingRule.CheckBitsEqualWidth.bad2.asl}


\RenderProseAndFormally{check_bits_equal_width}

\TypingRuleDef{PrecisionJoin}
\RenderRelation{precision_join}

\ExampleDef{Precision join}
In \listingref{precisionjoin}, the statement \verb|var b = (a * a) + 2;| is
forbidden because it tries to declare a type with a precision loss (see
\TypingRuleRef{LDVar}).
The expression \verb|a * a| has a type that results in a precision loss (see
\TypingRuleRef{AnnotateConstraintBinop}).
The typing rule \TypingRuleRef{PrecisionJoin} is called by
\TypingRuleRef{ApplyBinopTypes} to compute the precision of the type of the
expression \verb|(a * a) + 2|. Because the type of \verb|(a * a)| denotes a
precision lost, the type of \verb|(a * a) + 2| also denotes a precision lost.
\ASLListing{Precision join}{precisionjoin}{\typingtests/TypingRule.PrecisionJoin.asl}


\RenderProseAndFormally{precision_join}
